\documentclass[10pt,a4paper]{amsart}
\usepackage[margin=19mm]{geometry}
\usepackage{amsmath,amssymb,mathtools}
\usepackage[scr=rsfs,cal=euler]{mathalfa}
\usepackage{xfrac}
\usepackage[unicode,hidelinks,bookmarks=false]{hyperref}
\newcommand{\ie}{i.e.\ }
\newcommand{\cf}{cf.\ }
\newcommand{\resp}{respectively\ }
\newcommand{\Subsection}{\subsection}
\providecommand{\nobreakdash}{\nobreak\hbox{-}}
\setlength{\emergencystretch}{2em}
\multlinegap0pt
\usepackage{amssymb}
\usepackage{xcolor}
\usepackage{enumerate}
\newtheorem{proposition}{Proposition}
\newtheorem{lemma}{Lemma}
\newcommand{\AG}{\mathrm{AG}}
\newcommand{\PP}{\mathcal{P}}
\newcommand{\Q}{{\mathbb{Q}}}
\DeclareMathOperator{\dimQ}{dim_{\Q}}
\newcommand{\ideal}[1]{\left\langle #1 \right\rangle}
\newcommand{\lcm}{{\rm lcm}}
\newcommand{\lex}{lex}
\title{Gr\"obner Bases for Alexander Fitting Ideals of Links}
\author{Takefumi Nosaka, Kotaro Shimizu}
\date{Source: arXiv:2608.25725v1 (CC BY 4.0)}
\begin{document}
\maketitle

\noindent\emph{Source and licence.} Gr\"obner Bases for Alexander Fitting Ideals of Links. Version arXiv:2608.25725v1 (CC BY 4.0), distributed by arXiv under the Creative Commons Attribution 4.0 International licence, http://creativecommons.org/licenses/by/4.0/, as recorded in the arXiv licence metadata for this exact version and shown on the abstract page https://arxiv.org/abs/2608.25725v1. Full licence notice: \texttt{licenses/arxiv-2608.25725-CC-BY-4.0.txt}.\par\smallskip

\noindent\textit{Editorial scope.} Counted unit: the article's lemma lem:pretzel-local-uniform (source0.tex lines 1288--1330). The article's complete original proof of it is delivered with it (source0.tex lines 1332--1390, one \texttt{proof} environment of 59 lines). No mathematics is written, repaired or re-derived: the statement and the proof are the article's own lines, in the article's own notation, and no retained line is substituted or re-flowed.  Retained uncounted as the source-local context the proof needs: lines 1155--1169.  Nothing was omitted from any retained span, so the package carries no omission record.  No retained line contains a citation, so the wrapper carries no bibliography and no bibliography entry.  No retained line contains a figure environment, an image-inclusion command or drawing code, so no image is delivered and the package declares no asset dependency.  Editorial formatting only, in addition: an AMS \texttt{amsart} preamble that restates the article's own declarations and macros used by the retained lines, namely the environments \texttt{proposition}, \texttt{lemma} on separate counters, together with the standard packages those lines need.  The retained environments print in this document as Proposition 1, Lemma 1 in order of appearance, whereas the article prints the same items with the numbering of its own full document.  The complete native source of the article as distributed in the arXiv e-print for this exact version is bundled unchanged as \texttt{sources/208592/source0.tex}.
\begin{proposition}[Reduction to a three-generated ideal]
\label{prop:pretzel-reduction}
Let \(I:=\PP(E_2(L))\subset R\), where we use the identification
\(t_1=y\), \(t_2=x\).  Then
\[
I=\ideal{A_q(x),G_k(x,y),H_m(x,y)}.
\]
Consequently,
\[
\AG_2(L)=\mathcal G_{\lex(y>x)}(I),
\]
where \(\mathcal G_{\lex(y>x)}(I)\) denotes the reduced Gr\"obner basis of
\(I\subset R\) for the order \(\lex(y>x)\).  Moreover, in \(R/I\) one has
\((xy)^k=1\) and \(x^m=y^m\).
\end{proposition}

\begin{lemma}[Cyclotomic fibers]\label{lem:pretzel-local-uniform}
For every divisor \(\delta\mid q\) with \(\delta>1\),
\begin{equation}\label{eq:pretzel-local-factor}
p_\delta(y)=
\begin{cases}
1,
& \delta\nmid\lcm(m,k),
\\[6pt]
\displaystyle
\frac{y^g-x_\delta^e}
 {(y-x_\delta)^{\mathbf 1_{\delta\mid k}}
  (y-x_\delta^{-1})^{\mathbf 1_{\delta\mid m}}},
& \delta\mid\lcm(m,k).
\end{cases}
\end{equation}
In particular,
\begin{equation}\label{eq:pretzel-local-degree}
d_\delta=
\begin{cases}
0,
& \delta\nmid\lcm(m,k),
\\[3pt]
g-\mathbf 1_{\delta\mid m}-\mathbf 1_{\delta\mid k},
& \delta\mid\lcm(m,k).
\end{cases}
\end{equation}

Let
\[
\mathcal D:=
\{\delta>1:\,\delta\mid q,\ d_\delta>0\}.
\]
If \(\mathcal D=\varnothing\), then \(I=R\).  If
\(\mathcal D\ne\varnothing\), then
\[
R/I\cong
\prod_{\delta:\,\delta\in\mathcal D}
\mathcal K_\delta[y]/(p_\delta(y)),
\qquad
\dimQ(R/I)=
\sum_{\delta:\,\delta\in\mathcal D}d_\delta\deg\Phi_{2\delta}.
\]
\end{lemma}

\begin{proof}
Write \(\mu_n:=\{\zeta\in\overline{\Q}^{\,\times}:\,\zeta^n=1\}\).
The roots in \(y\) of the two specialized polynomials are simple and are
given by
\[
H_m(x_\delta,y)=0
\iff y=x_\delta\sigma
\quad(\sigma\in\mu_m\setminus\{1\}),
\]
\[
G_k(x_\delta,y)=0
\iff y=x_\delta^{-1}\tau
\quad(\tau\in\mu_k\setminus\{1\}).
\]
The simultaneous solutions of \(y^m=x_\delta^m\) and
\(y^k=x_\delta^{-k}\) are \(y=x_\delta\sigma\), where
\(\sigma\in\mu_m\cap x_\delta^{-2}\mu_k\).
To obtain the common roots of \(H_m(x_\delta,y)\) and
\(G_k(x_\delta,y)\), one must exclude
\(\sigma=1\) and \(\sigma=x_\delta^{-2}\), respectively.
This intersection is nonempty precisely when
\(x_\delta^{-2}\in\mu_m\mu_k=\mu_{\lcm(m,k)}\), or equivalently when
\(\delta\mid\lcm(m,k)\).  When nonempty, it is a
coset of \(\mu_m\cap\mu_k=\mu_g\), and hence has \(g\) elements.
For every \(\sigma\) in this coset ,the corresponding value \(y=x_\delta\sigma\) satisfies
\[
y^m=x_\delta^m,
\qquad
y^k=x_\delta^{-k},
\qquad
y^g=(y^m)^u(y^k)^v=x_\delta^e.
\]
Since \(y^g-x_\delta^e\) has exactly \(g\) simple roots, these \(g\)
values are precisely its roots.  The excluded root \(y=x_\delta\) occurs
precisely when \(\delta\mid k\), while the excluded root
\(y=x_\delta^{-1}\) occurs precisely when \(\delta\mid m\).  These roots
are distinct because \(\delta>1\).  This proves
\eqref{eq:pretzel-local-factor}, and
\eqref{eq:pretzel-local-degree} follows.

Since \(q\) is odd,
\[
A_q(x)=\prod_{\delta:\,\delta\mid q,\ \delta>1}\Phi_{2\delta}(x).
\]
The Chinese remainder theorem and
Proposition~\ref{prop:pretzel-reduction} therefore give
\[
R/I\cong
\prod_{\delta:\,\delta\mid q,\ \delta>1}
\frac{\mathcal K_\delta[y]}{(G_k(x_\delta,y),H_m(x_\delta,y))}.
\]
Because \(\mathcal K_\delta[y]\) is a PID, the ideal in the denominator
is \((p_\delta)\).  If \(\mathcal D=\varnothing\), every factor is the
zero ring, so \(I=R\).  Otherwise the zero factors, characterized by
\(p_\delta=1\), may be omitted.  Finally,
\(\dimQ(\mathcal K_\delta[y]/(p_\delta))
=d_\delta[\mathcal K_\delta:\Q]=d_\delta\deg\Phi_{2\delta}\), which
proves the dimension formula.
\end{proof}

\end{document}
